% Einheit 4 von 5 – Dezimalzahlen (bank/brueche-dezimalzahlen/e4.jsonl, zusammenbau v0.1)
%% TODO Zweigzeile Teil 1: was hier gelernt wird (Satz in Schülersprache aus dem Einheitstitel) – Einheit 4
\einheitenkopf[e4]{Einheit 4 von 5 · Dezimalzahlen}
\zweigzeile{ab Kl. 5 · P10 oft}
\setcounter{aufgabe}{20}

%% TODO Titel: Ich-kann-Satz fehlt in der Bank (Platzhalter: Kettenname) – Nr. 21
%% TODO Anweisung: ein Satz über der ersten Teilaufgabe fehlt in der Bank – Nr. 21
\begin{aufgabe}{Umwandeln}
\begin{teile}
\teil $\frac{9}{100}$ – wie viele Stellen nach dem Komma? \leerfeld
\end{teile}
%% TODO Darstellung für schwach fehlt in der Bank (brueche-dezimalzahlen-e4-k1-s1-v1; Abschnitt „Für schwache Schüler“)
\swz{$\frac{6}{10}$ – als Dezimalzahl?}{}
%% TODO Darstellung für schwach fehlt in der Bank (brueche-dezimalzahlen-e4-k1-s1-v2; Abschnitt „Für schwache Schüler“)
\swz{$\frac{47}{100}$ – als Dezimalzahl?}{}
%% TODO Darstellung für schwach fehlt in der Bank (brueche-dezimalzahlen-e4-k1-s1-v3; Abschnitt „Für schwache Schüler“)
\swz{$0{,}8$ – als Bruch mit dem Nenner $10$?}{}
%% TODO Darstellung für schwach fehlt in der Bank (brueche-dezimalzahlen-e4-k1-s1-v4; Abschnitt „Für schwache Schüler“)
\swz{$0{,}39$ – als Zehnerbruch?}{}
%% TODO Darstellung für schwach fehlt in der Bank (brueche-dezimalzahlen-e4-k1-s1-v5; Abschnitt „Für schwache Schüler“)
\swz{$\frac{83}{100}$ – als Dezimalzahl?}{}
\swfrage{Was bleibt gleich, was ändert sich?}
%% TODO Darstellung für schwach fehlt in der Bank (brueche-dezimalzahlen-e4-k1-s2-v1; Abschnitt „Für schwache Schüler“)
\swz{$\frac{7}{100}$ – als Dezimalzahl?}{}
\swa{Trage $0{,}4$ und $1{,}3$ am Zahlenstrahl ein.}{\zahlenstrahl[xmin=0,xmax=2,xstep=0.1,karo=0.6]{0/0, 1/1, 2/2}}
%% TODO Darstellung für schwach fehlt in der Bank (brueche-dezimalzahlen-e4-k1-s4-v1; Abschnitt „Für schwache Schüler“)
\swz{$\frac{2}{5}$ – als Dezimalzahl?}{}
%% TODO Darstellung für schwach fehlt in der Bank (brueche-dezimalzahlen-e4-k1-s5-v1; Abschnitt „Für schwache Schüler“)
\swz{$0{,}6$ – als gekürzter Bruch?}{}
%% TODO Darstellung für schwach fehlt in der Bank (brueche-dezimalzahlen-e4-k1-s6-v1; Abschnitt „Für schwache Schüler“)
\swz{$\frac{7}{8}$ – als Dezimalzahl? Teile den Zähler durch den Nenner.}{}
%% TODO Darstellung für schwach fehlt in der Bank (brueche-dezimalzahlen-e4-k1-s7-v1; Abschnitt „Für schwache Schüler“)
\swz{$\frac{4}{9}$ – als Dezimalzahl? Die Division geht nicht auf; schreibe mit Periodenstrich.}{}
\begin{teile}
\teil Welche Aussage ist wahr? Kreuze an. \hfill (P10 2015 OS)\\ \kreuz{$2{,}5 < \frac{5}{2}$}\\ \kreuz{$\frac{13}{5} > \frac{5}{2}$}\\ \kreuz{$\sqrt{5} > \frac{5}{2}$}
\end{teile}
\end{aufgabe}

%% TODO Titel: Ich-kann-Satz fehlt in der Bank (Platzhalter: Kettenname) – Nr. 22
%% TODO Anweisung: ein Satz über der ersten Teilaufgabe fehlt in der Bank – Nr. 22
\begin{aufgabe}{Zahl in die Stellenwerttafel eintragen und ablesen (Zehntel, Hundertstel, Tausendstel)}
\swa{Welche Dezimalzahl steht in der Stellenwerttafel? \leerfeld}{\sachtabelle{cccc}{Einer & Zehntel & Hundertstel & Tausendstel}{2 & 0 & 4 & 7}}
\end{aufgabe}

%% TODO Titel: Ich-kann-Satz fehlt in der Bank (Platzhalter: Kettenname) – Nr. 23
%% TODO Anweisung: ein Satz über der ersten Teilaufgabe fehlt in der Bank – Nr. 23
\begin{aufgabe}{Nachbarzahlen (Nachbar-Einer, -Zehntel, -Hundertstel) und Zählen in Schritten (0,2er-Schritte; über 2,9 hinaus)}
%% TODO Darstellung für schwach fehlt in der Bank (brueche-dezimalzahlen-e4-k3-s1-v1; Abschnitt „Für schwache Schüler“)
\swz{Nenne die Nachbarzehntel von $4{,}37$.}{}
\end{aufgabe}

%% TODO Titel: Ich-kann-Satz fehlt in der Bank (Platzhalter: Kettenname) – Nr. 24
%% TODO Anweisung: ein Satz über der ersten Teilaufgabe fehlt in der Bank – Nr. 24
\begin{aufgabe}{Umwandeln – Fehler finden, Begründen, Darstellungswechsel, Anwendung}
\swz[3]{Jana schreibt: $\frac{1}{4} = 1{,}4$. Finde den Fehler.}{}
\swfrage{Warum ist $\frac{9}{20} = 0{,}45$? Begründe.}
\swz{Welche Dezimalzahl zeigt der gefärbte Teil des Streifens?}{\bruchrechteck{2}{10}}
\swz[3]{Ein Rezept verlangt $\frac{3}{5}$ l Milch. Der Messbecher hat eine Skala in Litern mit Komma. Bis zu welcher Marke füllst du?}{}
\end{aufgabe}

\uebersichtskasten{Dezimalzahl: Die Stellen hinter dem Komma heißen Zehntel, Hundertstel, Tausendstel. \quad 0,7 = 7/10 \quad 0,05 = 5/100 \quad 1,25 = 1 25/100 = 1 1/4 \\ Bruch → Dezimalzahl: auf Nenner 10, 100 oder 1000 erweitern – oder Zähler durch Nenner teilen. \\ 3/20 = 15/100 = 0,15 \quad 5/8 = 5 : 8 = 0,625 \quad 2/3 = 2 : 3 = 0,666… (periodisch, Strich über der 6) \\ Merkzahlen: 1/2 = 0,5 \quad 1/4 = 0,25 \quad 3/4 = 0,75 \quad 1/5 = 0,2 \quad 1/8 = 0,125 \quad 1/3 = 0,333…}
